\centerline{\bf \Huge Пересечение фигур}

\problem{1} Какие геометрические фигуры могут образоваться при пересечении двух лучей?

\problem{2} Начертите в тетради два треугольника так, чтобы их пересечением были:
\textbf{а)} треугольник; \textbf{б)} прямоугольник; \textbf{в)} отрезок; \textbf{г)} точка.

\problem{3} Какие фигуры можно получить при пересечении двух треугольников?

\problem{4} Какие фигуры можно получить при пересечении двух четырёхугольников?

\problem{5} На рисунке два треугольника разделяют листок бумаги на 6 частей (шестая часть --- это то, что останется на листе, если вырезать оба треугольника). Нарисуйте два четырёхугольника, которые разделяют лист бумаги на 9 частей. Пронумеруйте полученные части.

\begin{center}
\includegraphics[width=0.2\textwidth]{Летние Сборы/Сборы Элиста 2024/ВВ/2022-09-23 треугольники.png}
\end{center}

\problem{6} Восемь бумажных квадратов $2 \times 2$ последовательно выкладывали на стол, пока не получился большой квадрат $4 \times 4$. Последним на стол положили квадрат $E$. На рисунке изображено, как видны квадраты: квадрат $E$ видно полностью, остальные квадраты видно частично. Какой квадрат положили на стол третьим по счёту?

    %не конструктив, но подходит
    \begin{center}
\includegraphics[width=0.2\textwidth]{Летние Сборы/Сборы Элиста 2024/ВВ/7 квадратов.png}
\end{center}
